\begin{claim}{D.2.44}{Proved}{The universal Moore composite comes from the bottom sphere}
Work at the prime two, write $P_n=M(\mathbb Z/2,n-1)$, and
let $i_n:S^{n-1}\to P_n$ be the bottom inclusion.  Put
\[
 w=[i_5,\operatorname{id}_{P_5}]:P_8\longrightarrow P_5,
 \qquad h_M=jw:P_8\longrightarrow Y=\Sigma^2K(\FF_2,2),
 \tag{D.2.44a}
\]
where $j:P_5\to Y$ is the bottom Moore map of D.2.41.
For the Stiefel attachment $l_4\in\pi_{14}P_8$ of that
claim, the following three assertions hold.

First, the actual universal composite satisfies
\[
 wl_4\in\operatorname{im}\bigl(
    i_{5*}:\pi_{14}S^4_{(2)}\longrightarrow\pi_{14}P_5\bigr).
 \tag{D.2.44b}
\]
Second, with $u_4=ji_5$ and $\alpha=u_4\nu_4$, its target
image has the form
\[
 h_Ml_4=a\,\alpha\sigma'+b\,u_4E\epsilon',
 \qquad a\in\mathbb Z,\quad b\in\{0,1\},
 \qquad 2u_4E\epsilon'=0.
 \tag{D.2.44c}
\]
Neither coefficient is evaluated here.  Third, the two bottom
images that remain in this formula already suspend to zero
separately:
\[
 E(i_5\nu_4\sigma')=0,
 \qquad E(i_5E\epsilon')=0
                 \quad\text{in }\pi_{15}P_6.
 \tag{D.2.44d}
\]
Consequently the known equality $E(wl_4)=0$ imposes no further
relation between these two coefficients.  These assertions do
not evaluate $h_Ml_4$ or decide the James--Hopf image question
of D.2.43.
\end{claim}
\begin{proof}
All Whitehead products below are actual products of spaces.
The generalized product in (D.2.44a) has domain
$\Sigma(S^3\wedge P_4)=P_8$, and its bottom restriction is
\[
 wi_8=i_5\omega_4,\qquad
 \omega_4=[\iota_4,\iota_4].
 \tag{D.2.44e}
\]
This is the generalized-Whitehead/commutator convention of
D.2.29 and D.2.41; see
\cite[Definitions 2.2--2.3 and Theorem 2.4,
pp.~10--13]{EH5Arkowitz1962}.
Morisugi--Mukai give
\[
 \pi_{14}P_8=\mathbb Z/4\{l_4\}
       \oplus\mathbb Z/2\{\widetilde{\nu_7^2}\},
 \quad 2l_4=i_8\sigma',\quad
 j_*l_4=[\gamma_8,\iota_7],
 \tag{D.2.44f}
\]
where the last $j_*$ is the map to the relative group
of $(P_8,S^7)$ and $\gamma_8$ is its characteristic
relative eight-cell
\cite[Lemmas 2.7--2.8, p.~523]{EH5MorisugiMukai2000}.
In particular the pinch $p_8:P_8\to S^8$ kills $l_4$:
in the relative product it kills the absolute bottom entry.

\medskip\noindent
\textbf{The actual third cell of the pinch fiber.}
Let $F_5$ be the homotopy fiber of $p_5:P_5\to S^5$.
Its cells begin in dimensions $4,8,12,16$, and the first
two nontrivial skeleta are
\[
 A_5=S^4\cup_{2\omega_4}e^8,
 \qquad B_5=A_5\cup_{2\tau_5}e^{12}.
 \tag{D.2.44g}
\]
The second formula is an exact two-local statement.
Indeed Mukai computes
\[
 \pi_{11}A_5=\mathbb Z\{\tau_5\}
   \oplus\mathbb Z/4\{\widetilde{2\nu_7}\}
   \oplus\mathbb Z/15\{i\phi\},
 \qquad j_*\tau_5=3[\gamma_8,\iota_4],
 \tag{D.2.44h}
\]
and proves that the actual twelve-cell attachment is
$2\tau_5+a i\phi$, with $\phi\in\pi_{11}S^4$ of
odd order.  In particular it has no component in the displayed
Moore coextension summand
\cite[Lemma 4.1(iii), p.~69, and Lemma 5.1(i), p.~70]
{FGHMukai2001}.  Thus (D.2.44g) does not discard a
two-primary indeterminacy in the attachment.

Choose the map of cofibers $c:P_8\to A_5$ induced by the
square with top map $2:S^7\to S^7$, bottom map
$2\omega_4:S^7\to S^4$, left vertical map the identity,
and right vertical map $\omega_4$.  Its bottom map is
$\omega_4$ and its relative top degree is one.  Write
$f:P_8\xrightarrow{c}A_5\to P_5$ for its actual composite.
The restrictions of $f$ and $w$ to $S^7$ agree, so exactness
of the Moore Puppe sequence gives
\[
                    w=f+h p_8
 \quad\text{for some }h:S^8\longrightarrow P_5.
 \tag{D.2.44i}
\]

\medskip\noindent
\textbf{The necessary Hopf--Hilton correction.}
Although $p_8l_4=0$, precomposition by $l_4$ need not
preserve addition in $[P_8,P_5]$.  We compute the correction
in (D.2.44i) before using it.  Write $P_8=\Sigma P_7$
and apply Hilton--Milnor to
$\Omega\Sigma(P_7\vee P_7)$ through degree fourteen.
Besides the two singleton factors, the only possible factor
is $\Sigma(P_7\wedge P_7)$, with bottom cell thirteen.
Word length three begins in dimension nineteen and cannot
contribute.  This is the usual two-letter Hilton--Milnor
decomposition, with the range obtained from the indicated
bottom dimensions; compare
\cite[Proposition 2.2 and its proof]{I_CMN1987}.

Since the second map is $h p_8$, the mixed factor further
maps through
\[
       \Sigma(P_7\wedge S^7)\simeq P_{15}.
 \tag{D.2.44j}
\]
Its fourteenth homotopy group is its bottom $\mathbb Z/2$
by Hurewicz.  Naturality of the Whitehead map therefore
identifies the entire possible correction with
$\varepsilon[\alpha_0,h]$, where
$\alpha_0=i_5\omega_4=[i_5,i_5]$ and
$\varepsilon\in\FF_2$.  There is no omitted top--top
or higher-word contribution in this range.

Put $u=i_5$.  Its shifted degree is three, hence odd.
Graded Jacobi and bilinearity give
\[
 [[u,u],h]=\pm2[u,[u,h]]
          =\pm[2u,[u,h]]=0,
 \tag{D.2.44k}
\]
since $2u=0$ in the actual bottom Moore homotopy group.
Thus the correction is zero and the required equality is
valid:
\[
                 wl_4=(A_5\longrightarrow P_5)cl_4.
 \tag{D.2.44l}
\]
This argument is essential: the pinch equation alone would
not justify (D.2.44l).

\medskip\noindent
\textbf{Relative Jacobi and the bottom image.}
Naturality of the relative Whitehead product gives
\[
 j_*(cl_4)=[\gamma_8,\omega_4]
          =2[[\gamma_8,\iota_4],\iota_4]
 \tag{D.2.44m}
\]
up to the common orientation convention.  Only one relative
class occurs here.  The mixed relative Whitehead module
satisfies the Jacobi derivation law; the shifted degrees of
$\gamma_8$ and $\iota_4$ are seven and three, respectively,
so its two terms have the same sign.  James proves precisely
this mixed relative Samelson identity, and its equivalent
left-action form, in
\cite[\S4, equations (4.1)--(4.2), p.~337]{FGHJames1971}.
The loop-group model of a CW pair identifies it with the
relative Whitehead module identity used here.  No product
of two arbitrary relative classes or metastable exact
sequence is being assumed.

Combining (D.2.44h) and (D.2.44m) yields
\[
           3j_*(cl_4)=j_*[2\tau_5,\iota_4].
 \tag{D.2.44n}
\]
In $B_5$ the absolute class $2\tau_5$ is zero by the
actual twelve-cell attachment.  Therefore the image of
$3j_*(cl_4)$ in $\pi_{14}(B_5,S^4)$ is zero.  Multiplication
by three is invertible at the prime two.  Exactness of this
pair puts the image of $cl_4$ in $B_5$ in the bottom-sphere
image.  Composing $B_5\to F_5\to P_5$ and using
(D.2.44l) proves (D.2.44b).

\medskip\noindent
\textbf{The two remaining target terms.}
The sphere table and composition relation needed here are
\[
\begin{split}
 \pi_{14}S^4_{(2)}
    &=\mathbb Z/8\{\nu_4\sigma'\}
      \oplus\mathbb Z/4\{E\epsilon'\}
      \oplus\mathbb Z/2\{\eta_4\mu_5\},\\
 (E\nu')\sigma'&=2E\epsilon'.
\end{split}
 \tag{D.2.44o}
\]
They are displayed and proved, respectively, in
\cite[Propositions 7 and 13, pp.~7--9]{FGHStricklandToda}.
For clarity about the source's Hopf-formula warning, the
left factor $E\nu'$ is a suspension.  For the comparison
$H(\nu_4\sigma')=\sigma'$ one may use the formula after
one suspension and the injectivity of
$E:\pi_{14}S^7\to\pi_{15}S^8$, since both $\sigma'$
and $E\sigma'$ have order eight.  The remaining correction
$\nu_7\nu_{10}\eta_{13}$ is zero.  Thus this instance
does not use the unqualified general composite formula
about which that source warns.

The established target relations are
$u_4\eta_4=0$, $\alpha=u_4\nu_4$, $4\alpha=0$, and
$[u_4,u_4]=-2\alpha$ (D.2.17--D.2.18 and D.2.29).
Using $\omega_4=2\nu_4-E\nu'$ with compatible signs gives
\[
 u_4E\nu'=2\alpha-[u_4,u_4]=4\alpha=0,
 \qquad
 2u_4E\epsilon'=u_4(E\nu')\sigma'=0.
 \tag{D.2.44p}
\]
Postcomposition by $u_4$ is a homomorphism on the sphere
homotopy group.  It kills the third summand in (D.2.44o),
and (D.2.44p) reduces the coefficient of the second to
zero or one.  The bottom-image assertion now proves
(D.2.44c).  The exact double from D.2.41 remains
$2(h_Ml_4)=(-2\alpha)\sigma'$; it is not an evaluation
of either coefficient.

\medskip\noindent
\textbf{Why suspension supplies no additional coefficient test.}
The same primary table gives
\[
 E(\nu_4\sigma')=\pm E^2\epsilon'
                 =2\nu_5\sigma_8.
 \tag{D.2.44q}
\]
We show that the bottom inclusion $i_6$ kills this common
class.  Since $\nu_5$ is a suspension, the degree-two map
commutes with it.  The Hopf invariant of $\sigma_8$ is one,
so the actual Hopf--Hilton formula gives
\[
 (2\iota_5)\nu_5\sigma_8
 =\nu_5(2\iota_8)\sigma_8
 =2\nu_5\sigma_8+[\nu_5,\nu_5].
 \tag{D.2.44r}
\]
There are no higher Hilton terms in the dimension of
$\sigma_8$.  Naturality and the classical identity
$[\iota_5,\iota_5]=\nu_5\eta_8$ imply
\[
 [\nu_5,\nu_5]
 =\pm[\iota_5,\iota_5]\,\Sigma(\nu_4\wedge\nu_4)
 =\pm\nu_5\eta_8\nu_9\nu_{12}=0.
 \tag{D.2.44s}
\]
Here $\eta_8\nu_9=0$ because it belongs to the stable
fourth stem, which is zero.  The Whitehead identity is
recorded, with its exact Toda equations (5.10) attribution,
in the accompanying author program of
\cite{FGHStricklandToda}.  Finally $i_6(2\iota_5)=0$
by the Moore attachment, so (D.2.44r)--(D.2.44s) give
$i_6(2\nu_5\sigma_8)=0$.  Equation (D.2.44q) proves
(D.2.44d).

Thus the suspension test vanishes on each of the two
candidate bottom images separately.  In particular it
cannot replace $E(i_5\delta)=0$ by the stronger assertion
$E\delta=0$ for a bottom representative $\delta$.
The actual coefficient problem and the full target
James--Hopf image remain as stated in D.2.43.
\end{proof}
